% =============================================================================
% SOICT 2026 Submission — Springer CCIS / LNCS Format
% Track: AI Foundations, Foundation Models, and Generative AI
% Format: Single-Blind Review (Author names & affiliations included)
% Page Limit: Max 12 pages excluding references
% Template: \documentclass[runningheads]{llncs}
% =============================================================================
\documentclass[runningheads]{llncs}

\usepackage[utf8]{inputenc}
\usepackage[T5,T1]{fontenc}
\usepackage[vietnamese,english]{babel}
\usepackage{listings}
\usepackage{amsmath,amssymb,amsfonts}
\usepackage{graphicx}
\usepackage{booktabs}
\usepackage{multirow}
\usepackage[hyphens]{url}
\usepackage{microtype}
\usepackage[hidelinks]{hyperref}

\newcommand{\vn}[1]{{\fontencoding{T5}\selectfont #1}}

\begin{document}
\pagestyle{empty}

\title{Early-Stopping Activation Steering for Factual Preference Alignment in Vietnamese Medical Language Models}
\titlerunning{Early-Stopping Activation Steering for Vietnamese Medical LMs}

\author{Phan Do Thanh Tuan}
\authorrunning{P. D. T. Tuan}

\institute{Department of Information Technology, FPT University Quy Nhon, Binh Dinh, Vietnam \\
\email{phandothanhtuan1704@gmail.com}}

\maketitle

\begin{abstract}
Small Language Models (SLMs) offer practical utility for clinical guidance in low-resource settings, but frequently exhibit domain-specific hallucinations in non-English contexts like Vietnamese. We propose Early-Stopping Activation Steering, an inference-time framework evaluating temporally bounded vector injection at Layer~8 of Qwen2.5-7B-Instruct under Hard Cutoff and Linear Decay schedules ($K{=}16$, decoding steps $t \in [1, 16]$). Under natural completion ($T=800$), Hard Cutoff ($K{=}16$) achieves 70.60\% automated reference-preference accuracy (+5.20~pp over baseline 65.40\%, item-level exact McNemar $p=0.00086$, Holm $p_{\text{adj}}=0.0026$) with 100\% EOS hit rate; however, longer outputs coincide with lower token-level BERTScore F1 (0.6695 vs 0.6779) and higher 4-gram repetition (5.35\% vs 4.10\%), a divergence whose causal mechanism remains unestablished. Under bounded stress testing ($T=200$), Linear Decay ($K{=}16$) achieves 77.40\% automated reference preference (+4.20~pp over baseline 73.20\%, item-level exact McNemar $p=0.0055$, Holm $p_{\text{adj}}=0.0165$, 95\% bootstrap CI $[+1.40, +7.00]$~pp), showing no statistically significant separation from multi-retriever Hybrid RRF RAG (76.20\%, $p=0.4709$) and exceeding lexical BM25 RAG (+8.80~pp, $p=3.85 \times 10^{-7}$) in an operational pipeline comparison. An exploratory pipeline combining Hybrid RAG with Hard Cutoff yields 80.40\% reference preference (+1.00~pp over RAG alone, $p=0.3018$, non-significant marginal gain; 4.46\% Rep-4 vs 4.63\%). Directional controls establish a standardized descriptive separation of $z_{\text{sep}}=5.20\sigma$ over isotropic noise ($55.82\% \pm 4.15\%$). While early-stopping steering provides a lightweight, retrieval-free preference alignment mechanism, evaluations reflect the joint effect of schedule and cumulative dose within a single model family, warranting cross-architecture and clinical validation.

\keywords{Activation steering \and Automated reference preference \and Clinical language models \and Vietnamese NLP \and Inference-time intervention.}
\end{abstract}

\section{Introduction}\label{sec:intro}
Deploying Large Language Models (LLMs) in high-stakes clinical domains demands rigorous factual fidelity~\cite{ref3,ref13}. When adapted to low-resource languages such as Vietnamese, models frequently generate plausible yet clinically hazardous hallucinations. While parameter-efficient fine-tuning (PEFT)~\cite{ref7} and retrieval-augmented generation (RAG)~\cite{ref5} are standard countermeasures, they exhibit operational trade-offs: Supervised Fine-Tuning (SFT)~\cite{ref10} risks catastrophic forgetting of general reasoning capabilities~\cite{ref12}, while RAG introduces substantial retrieval latency (118--142~ms in our benchmarks) and external retrieval infrastructure overhead~\cite{ref_faiss}.

Representation Engineering (RepE)~\cite{ref1}, Activation Addition~\cite{ref8}, and Inference-Time Intervention (ITI)~\cite{ref2,ref9} offer a lightweight alternative: steering internal hidden representations during inference without updating model parameters~\cite{ref1,ref2}, leveraging internal hidden representations that correlate with factual truthfulness~\cite{ref11}. However, conventional continuous steering (applying an intervention vector across all generated tokens) often leads to semantic drift, cognitive incoherence, and repetitive output loops.

To address these limitations, we investigate \textbf{Early-Stopping Activation Steering}. We hypothesize that factual commitments are established during initial decoding steps; consequently, applying activation steering exclusively during an early temporal window ($K=16$ tokens) should be sufficient to steer generation toward factual paths while allowing subsequent decoding to terminate naturally. We ground our evaluation on \textbf{ViHaluEval-Medical}, a curated benchmark of 14,674 factual-counterfactual question-answer pairs derived from the official Vietnamese National Drug Formulary (\vn{Dược thư Quốc gia Việt Nam}, 2nd Edition, 2018)~\cite{ref6}. We target Qwen2.5-7B-Instruct~\cite{ref4} as an edge-deployable Small Language Model (SLM) suitable for single-GPU deployment on an enterprise accelerator (NVIDIA Tesla T4 16~GB, with 4-bit NF4 weights occupying 4.86~GB allocated and 5.52~GB peak reserved VRAM during $T=800$ generation~\cite{ref_bnb}).

\noindent\textbf{Key Contributions:}
\begin{enumerate}
    \item \textbf{Curated Vietnamese Medical Benchmark:} We construct ViHaluEval-Medical (14,674 rule-filtered contrastive pairs from 14,700 candidates across 3 clinical categories). Manual formulary-grounding audit on $N=50$ source pairs yields 94.0\% alignment (47/50, 95\% Clopper-Pearson CI $[83.5\%, 98.7\%]$). Evaluations are drawn from a primary-API-disjoint held-out partition of 2,179 pairs, from which a balanced 500-item evaluation sample was drawn ($N_{\text{test}}=500$: 175 Special Dosage, 165 Pregnancy, 160 Drug Interactions across 373 distinct question clusters).
    \item \textbf{Early-Stopping Steering Dynamics:} We introduce temporally bounded intervention schedules (Hard Cutoff and Linear Decay, $K=16$). Under natural completion ($T=800$), Hard Cutoff achieves 70.60\% automated reference preference (+5.20~pp over baseline 65.40\%, item-level exact McNemar $p=0.00086$, Holm $p_{\text{adj}}=0.0026$). Under bounded stress testing ($T=200$), Linear Decay achieves 77.40\% (+4.20~pp over baseline 73.20\%, item-level exact McNemar $p=0.0055$, Holm $p_{\text{adj}}=0.0165$).
    \item \textbf{Directional Selectivity \& Operational Comparison:} Directional controls ($N=100$ isotropic vectors, $55.82\% \pm 4.15\%$) establish a standardized descriptive separation of $z_{\text{sep}}=5.20\sigma$ ($p_{\text{emp}}=0.0099$). In an operational pipeline comparison, steering matches multi-retriever Hybrid RRF RAG (77.40\% vs 76.20\%, $p=0.4709$) while eliminating external retrieval indexing and prompt expansion.
\end{enumerate}

\section{Methodology}\label{sec:methodology}
\subsection{Contrastive Activation Vector Estimation}
Let $\mathcal{D}_{\text{train}} = \{(x_i, y_i^+, y_i^-)\}_{i=1}^{N_{\text{train}}}$ denote contrastive training triplets ($N_{\text{train}}=10,272$), where $x_i$ is a clinical prompt, $y_i^+$ is a factual reference answer, and $y_i^-$ is a synthetic counterfactual hallucination generated following contrastive falsehood elicitation principles~\cite{ref16}. Given an autoregressive language model with $L_m=28$ transformer layers indexed by zero-based block index $j \in \{0, \dots, L_m-1\}$, let $h_j(x_i, y_i) = \operatorname{LayerOutput}_j(x_i \oplus y_i)[:, -1, :]$ denote the hidden state at residual block $j$ extracted at the sequence-terminal token position ($[:, -1, :]$, corresponding to the ChatML sequence terminator \texttt{<|im\_end|>}). The contrastive difference vector at block $j$ is:
\begin{equation}
\Delta h_j^{(i)} = h_j(x_i, y_i^+) - h_j(x_i, y_i^-) \in \mathbb{R}^d.
\end{equation}
The raw steering direction $\mu_j$ is the sample mean of contrastive differences:
\begin{equation}
\mu_j = \frac{1}{N_{\text{train}}} \sum_{i=1}^{N_{\text{train}}} \Delta h_j^{(i)}, \qquad v_{\text{steer}}^{(j)} = \frac{\mu_j}{\|\mu_j\|_2}.
\end{equation}
Layer-wise linear probing on the validation split ($N_{\text{val}}=2,223$, binary logistic regression classifying factual vs counterfactual terminal states across candidate blocks $j \in \{0, \dots, 27\}$) identified block $j^*=8$ (the 9th residual block, $d=3,584$) as exhibiting peak separation (84.6\% probe accuracy). Steering scale ($\alpha_0=18.0$) and horizon ($K=16$) were selected via validation-split grid search ($K \in \{8, 16, 24, 32\}$, $\alpha_0 \in \{10.0, 15.0, 18.0, 20.0\}$) prior to frozen test evaluation.

\subsection{Temporally Bounded Intervention Schedules}
During autoregressive generation, let $t \in \{1, 2, \dots, L_i\}$ denote the post-prefill decoding step index, where $L_i \in [1, T_{\text{gen}}]$ represents the completion length in generated decoding steps for prompt $i$. The prompt prefill forward pass remains unsteered, predicting initial token $y_1$; subsequent decoding steps $t \ge 1$ steer residual states at block $j^*=8$, shaping output predictions for tokens $y_2, \dots, y_{L_i+1}$ ($y_2, \dots, y_{17}$ under $K=16$). At layer $j^*=8$, the modified hidden state $\hat{h}_{j^*}^{(t)}$ is computed as:
\begin{equation}
\hat{h}_{j^*}^{(t)} = h_{j^*}^{(t)} + \alpha(t) \cdot v_{\text{steer}}^{(j^*)},
\end{equation}
where $\alpha(t)$ is governed by one of three temporal intervention schedules parameterized by initial strength $\alpha_0 \in \mathbb{R}^+$ (with finite-window schedules parameterized by cutoff horizon $K \in \mathbb{N}^+$):
\begin{align}
\textbf{Continuous Steering:} \quad &\alpha_{\text{cont}}(t) = \alpha_0, \quad \forall t \ge 1, \\
\textbf{Hard Cutoff Steering ($K \in \mathbb{N}^+$):} \quad &\alpha_{\text{hard}}(t) = \begin{cases} \alpha_0, & 1 \le t \le K, \\ 0, & t > K, \end{cases} \\
\textbf{Linear Decay Steering ($K \in \mathbb{N}^+$):} \quad &\alpha_{\text{decay}}(t) = \begin{cases} \alpha_0 \left(1 - \frac{t-1}{K}\right), & 1 \le t \le K, \\ 0, & t > K. \end{cases}
\end{align}
For finite-window schedules ($K \in \mathbb{N}^+$), $\alpha(t) = 0$ for $t > K$, leaving subsequent forward passes unmodified while preserving conditioned KV-cache states~\cite{ref_vllm}, whereas continuous steering bypasses $K$ and injects $\alpha_0$ across all decoding steps.

\begin{figure}[!t]
\centering
\begin{lstlisting}[language=Python,basicstyle=\ttfamily\scriptsize,frame=single,breaklines=true,columns=fullflexible]
class EarlyStoppingSteeringHook:
    def __init__(self, v_steer: torch.Tensor, alpha_0: float = 18.0, K: int = 16, schedule: str = "linear_decay"):
        self.v_steer = v_steer.to(dtype=torch.bfloat16)
        self.alpha_0, self.K, self.schedule = alpha_0, K, schedule
        self.step = 0  # reset per completion before generation

    def reset(self):
        self.step = 0

    def __call__(self, module, inputs, output):
        h = output[0] if isinstance(output, tuple) else output
        if h.shape[1] > 1:  # prefill forward pass remains unsteered
            return output
        self.step += 1
        if self.schedule != "continuous" and self.step > self.K:
            return output
        coeff = self.alpha_0 if self.schedule in ("hard_cutoff", "continuous") else self.alpha_0 * (1.0 - (self.step - 1) / self.K)
        h[:, -1, :] = h[:, -1, :] + coeff * self.v_steer.to(h.device)
        return (h,) + output[1:] if isinstance(output, tuple) else h
\end{lstlisting}
\caption{PyTorch implementation of the early-stopping activation steering forward hook (caller invokes \texttt{hook.reset()} before each generation pass).}
\label{lst:hook}
\end{figure}

\section{Experimental Setup}\label{sec:setup}
\subsection{ViHaluEval-Medical Benchmark Construction}
ViHaluEval-Medical was constructed from the Vietnamese National Drug Formulary (2nd Edition, 2018). A raw candidate pool of 14,700 QA pairs was extracted across 3 clinical categories (5,010 Special Dosage, 4,893 Pregnancy Safety, and 4,797 Drug Interactions). A rule-based filtering pipeline eliminated 26 malformed records (14 regex formatting anomalies, 8 duplicate pairs, 4 corrupted font encodings), yielding 14,674 rule-filtered records: Special Dosage (5,000 pairs, 34.07\%), Pregnancy \& Lactation Safety (4,885 pairs, 33.29\%), and Clinically Significant Drug Interactions (4,789 pairs, 32.64\%). The dataset comprises a 12,495-item development split (85.15\%, $N_{\text{train}}=10{,}272$, $N_{\text{val}}=2{,}223$) and a 2,179-item primary-API-disjoint held-out partition ($N_{\text{test\_partition}}=2{,}179$, 14.85\%). From this partition, a balanced evaluation sample of $N_{\text{test}}=500$ pairs was drawn (175 Special Dosage, 165 Pregnancy Safety, 160 Drug Interactions across 373 distinct question clusters). A manual formulary-grounding audit of $N=50$ randomly sampled training pairs directly against official monograph records confirmed 47/50 correct alignments (94.0\% alignment, 95\% Clopper-Pearson exact CI $[83.5\%, 98.7\%]$). All code, logs, and reproducible scripts are open-sourced at \url{https://github.com/dothanhtuan1704-droid/early-stopping-steering}.

\subsection{Evaluation Metrics}
\textbf{1.~Automated BERTScore Reference Preference (RefPref):} We compute token-level cosine similarities between generated output $\hat{y}_i$ and reference $y_i^+$ versus hallucination $y_i^-$ using multilingual BERT (`bert-base-multilingual-cased`, Layer~9)~\cite{ref15}:
\begin{equation}
\operatorname{RefPref}(\hat{y}_i) = \mathbb{I}\left[\operatorname{BERTScore}_{\text{F1}}(\hat{y}_i, y_i^+) > \operatorname{BERTScore}_{\text{F1}}(\hat{y}_i, y_i^-)\right].
\end{equation}
Aggregate accuracy is $\frac{1}{N} \sum_{i=1}^N \operatorname{RefPref}(\hat{y}_i) \times 100\%$. RefPref provides an automated embedding-similarity proxy against synthetic references, without direct clinician adjudication of numerical doses or contraindications.

\noindent\textbf{2.~Lexical \& Repetition Metrics:} We report ROUGE-L F1~\cite{ref14}, 4-gram repetition rate ($\text{Rep-4} = 1 - \frac{|\text{unique 4-grams}|}{|\text{total 4-grams}|}$), and EOS Hit Rate (percentage of completions emitting a valid end-of-sequence token).

\noindent\textbf{3.~Statistical Hypothesis Testing:} Statistical significance is assessed via paired exact McNemar tests on $2\times 2$ contingency tables ($a$: mutual win, $b$: baseline win, $c$: steered win, $d$: mutual loss). Multi-comparison error rates within evaluation panels are controlled via Holm-Bonferroni step-down correction ($p_{\text{adj}}$). Paired 95\% confidence intervals are estimated via $B=10,000$ bootstrap resamples, representing exploratory item-level inference.

\section{Results and Empirical Analysis}\label{sec:results}
\subsection{Primary Evaluation: Bounded vs High-Cap Horizons}
Table~\ref{tab:main_results} reports primary evaluations across two token horizons under greedy decoding: bounded stress testing ($T=200$) and natural completion ($T=800$).

\begin{table}[!t]
\caption{Evaluation on ViHaluEval-Medical ($N_{\text{test}}=500$) across intervention schedules.}
\label{tab:main_results}
\centering
\small
\resizebox{\textwidth}{!}{%
\begin{tabular}{lcccccc}
\toprule
\textbf{Intervention Schedule} & \textbf{RefPref (\%)} & \textbf{BERT F1} & \textbf{ROUGE-L / EOS Hit (\%)} & \textbf{Rep-4 (\%)} & \textbf{McNemar $p$} & \textbf{95\% CI (pp)} \\
\midrule
\multicolumn{7}{c}{\textit{\textbf{Panel A: Bounded Horizon Stress Test ($T=200$)}}} \\
\midrule
Unsteered Baseline & 73.20 & 0.7134 & 23.12 & \textbf{3.67} & -- & -- \\
Continuous ($K{=}\infty$) & 75.60 & 0.7133 & \textbf{23.36} & 3.90 & 0.1038 & $[-0.20, +5.00]$ \\
Hard Cutoff ($K{=}16$) & 77.00 & 0.7151 & 23.16 & 3.89 & 0.0110 & $[+1.00, +6.60]$ \\
\textbf{Linear Decay ($K{=}16$)} & \textbf{77.40} & \textbf{0.7157} & 23.33 & 3.88 & \textbf{0.0055} & $\mathbf{[+1.40, +7.00]}$ \\
\midrule
\multicolumn{7}{c}{\textit{\textbf{Panel B: High-Cap Natural Completion ($T=800$)}}} \\
\midrule
Unsteered Baseline & 65.40 & \textbf{0.6779} & \textbf{100.00} & \textbf{4.10} & -- & -- \\
Continuous ($K{=}\infty$) & 68.60 & 0.6712 & 99.80 & 4.38 & 0.0365 & $[+0.40, +6.00]$ \\
Linear Decay ($K{=}16$) & 67.80 & 0.6732 & 99.80 & 5.37 & 0.1189 & $[-0.20, +5.00]$ \\
\textbf{Hard Cutoff ($K{=}16$)} & \textbf{70.60} & 0.6695 & \textbf{100.00} & 5.35 & \textbf{0.00086} & $\mathbf{[+2.40, +8.20]}$ \\
\bottomrule
\end{tabular}%
}
\smallskip\noindent\scriptsize Note: Evaluated on disjoint test split $N_{\text{test}}=500$. Item-level paired 2$\times$2 counts ($a,b,c,d$ where $b$: baseline win, $c$: steered win): Panel A Linear Decay (350, 16, 37, 97, $p_{\text{adj}}=0.0165$), Hard Cutoff (350, 16, 35, 99, $p_{\text{adj}}=0.0219$), Continuous (349, 17, 29, 105, $p_{\text{adj}}=0.1038$). Panel B Hard Cutoff (311, 16, 42, 131, $p_{\text{adj}}=0.0026$), Continuous (309, 18, 34, 139, $p_{\text{adj}}=0.0730$), Linear Decay (308, 19, 31, 142, $p_{\text{adj}}=0.1189$). Holm-adjusted item-level $p_{\text{adj}}$ values reported; item-level 95\% CIs via $B=10,000$ paired bootstrap resamples. Col.~4: ROUGE-L (Panel A), EOS Hit Rate (Panel B).
\end{table}

Under bounded stress testing (Panel A, $T=200$), Linear Decay steering ($K{=}16$) achieves peak performance at \textbf{77.40\% RefPref} (+4.20~pp over baseline 73.20\%, item-level exact McNemar $p=0.0055$, Holm $p_{\text{adj}}=0.0165$, 95\% bootstrap CI $[+1.40, +7.00]$~pp). Under natural completion (Panel B, $T=800$), models maintain high EOS Hit rates (99.80\%--100.00\%): Hard Cutoff ($K{=}16$) achieves peak observed accuracy at \textbf{70.60\% RefPref} (+5.20~pp over baseline 65.40\%, item-level exact McNemar $p=0.00086$, Holm $p_{\text{adj}}=0.0026$, 95\% bootstrap CI $[+2.40, +8.20]$~pp), whereas Linear Decay yields 67.80\% and continuous steering yields 68.60\% ($p=0.0365$, Holm $p_{\text{adj}}=0.0730$). Notably, while pairwise RefPref preference increases by +5.20~pp, absolute BERTScore F1 against reference answers $y^+$ slightly decreases (0.6695 vs 0.6779) alongside a modest rise in Rep-4 (5.35\% vs 4.10\%). We hypothesize that this divergence relates to output length dynamics (mean 325.6 vs 248.7 Qwen subwords); however, because token-level Precision and Recall were not separately logged, attributing the F1 drop to output length remains an unverified hypothesis rather than an established mechanism.

\subsection{Empirical Activation-Norm Dynamics}\label{sec:norm_dynamics}
To evaluate activation mechanics at Layer~8, we assess hidden norms across $N=50$ diagnostic prompts ($T=100$):
Under unsteered baseline ($\alpha=0$), whole-phase mean magnitude is $\bar{H}_{8,\text{active}} = 52.88$ ($t \in [1, 16]$) and $\bar{H}_{8,\text{post16}} = 52.20$ ($t \in [17, L_i]$). Continuous steering ($K=\infty$) induces persistent elevation ($\Delta_{\text{norm}} = +1.90$ active, $+2.51$ post-16). Step 20 injection under extended schedules yields corresponding offsets ($\Delta_{\text{norm}}^{(20)}=+1.12$ for $K=24, \alpha=3.75$; $+1.98$ for $K=32, \alpha=7.31$). Post-cutoff ($t > K$), $\alpha(t)=0$ yields $\Delta_{\text{norm}} = 0.00$ by construction (as direct vector injection stops by definition), while downstream prefix conditioning persists in the KV-cache.

Evaluating prompt-level pre-hook norms under Hard Cutoff against unsteered baseline completions across $N=50$ paired prompt completions over post-cutoff steps $t \in [17, L_i]$ reveals prompt-level post-cutoff pre-hook norms under Hard Cutoff ($51.84$) exhibit no statistically significant difference relative to baseline ($52.20$, mean difference $-0.36$, sample SD $= 1.54$, paired $t$-test $p = 0.1050 > 0.05$, $95\%$ bootstrap CI $[-0.80, +0.08]$). While $p > 0.05$ fails to reject equal mean scalar norm magnitudes rather than establishing dynamic trajectory equivalence, the comparison demonstrates no statistically detectable difference in mean scalar activation norms ($51.84$ vs $52.20$, $p=0.1050$).

\subsection{Directional Controls \& Multi-Retriever RAG Comparisons}\label{sec:baselines_rag}
Table~\ref{tab:baselines} reports performance across controls and multi-retriever RAG. Directional controls were evaluated under matched $T=200, K=16, \alpha_0=18.0$ using identical Linear Decay injection. Negative steering reduces RefPref to 62.80\% ($-10.40$~pp vs baseline). Isotropic controls ($N=100$, $55.82\% \pm 4.15\%$ sample SD) establish standardized descriptive directional separation ($z_{\text{sep}}=5.20\sigma$, descriptive empirical rank $1/101$, $p_{\text{emp}} = 0.0099$). Label-shuffled steering ($v_{\text{shuf}}$, 68.60\%) indicates sensitivity to contrastive pair semantics. Furthermore, covariance-matched vectors ($N=20$, $74.81\% \pm 0.95\%$ sample SD) achieve +1.61~pp over baseline despite near-zero mean cosine similarity with $v_{\text{steer}}$ ($-0.06$). Because zero-mean Gaussian covariance sampling is centrally symmetric ($v \leftrightarrow -v$), sampling from $\Sigma_{\text{data}}$ restricts perturbations to the empirical subspace of contrastive data variance (avoiding isotropic collapse), while the oriented centroid $+v_{\text{steer}}$ ($77.40\%$) provides directed alignment (+2.59~pp / $2.73\sigma$ above the covariance mean), though $N=20$ samples provide a bounded preliminary baseline.

\begin{table}[!t]
\caption{Comparison against directional controls, placebo baselines, and multi-retriever RAG ($N_{\text{test}}=500$, $T=200$).}
\label{tab:baselines}
\centering
\small
\resizebox{\textwidth}{!}{%
\begin{tabular}{lcccccc}
\toprule
\textbf{Method / Condition} & \textbf{Raw Count} & \textbf{RefPref (\%)} & \textbf{Recall@1} & \textbf{MRR@10} & \textbf{BERT F1} & \textbf{Latency (s)$^\dagger$} \\
\midrule
Isotropic Gaussian ($N{=}100$) & -- & $55.82\pm4.15$ & -- & -- & 0.6945 & $23.33\pm0.30$ \\
Negative Steered ($-18.0 \cdot v_{\text{steer}}$) & 314 / 500 & 62.80 & -- & -- & 0.6889 & $23.34\pm0.28$ \\
BM25 Lexical RAG & 343 / 500 & 68.60 & 19.20\% & 0.2433 & 0.6970 & $7.32\pm0.45$ \\
Label-Shuffled ($v_{\text{shuf}}$) & 343 / 500 & 68.60 & -- & -- & 0.6845 & $23.34\pm0.28$ \\
Unsteered Baseline & 366 / 500 & 73.20 & -- & -- & 0.7134 & $23.33\pm0.30$ \\
Cov-Matched ($N{=}20$) & -- & $74.81\pm0.95$ & -- & -- & 0.7182 & $23.34\pm0.28$ \\
Dense BGE-M3 RAG~\cite{ref_bgem3} & 374 / 500 & 74.80 & 42.60\% & 0.5124 & 0.7180 & \textbf{7.27$\pm$0.42} \\
Hybrid RRF RAG & 381 / 500 & 76.20 & \textbf{48.20\%} & \textbf{0.5681} & \textbf{0.7190} & $7.39\pm0.48$ \\
\textbf{Linear Decay Steered} & \textbf{387 / 500} & \textbf{77.40} & -- & -- & 0.7157 & $23.34\pm0.28$ \\
Oracle Gold-Context RAG & 447 / 500 & 89.40 & 100.0\% & 1.0000 & 0.8002 & $7.17\pm0.38$ \\
\bottomrule
\end{tabular}%
}
\smallskip\noindent\scriptsize Note: Evaluated on disjoint test split $N_{\text{test}}=500$. Values with $\pm$ denote sample mean $\pm$ sample SD. Negative steering injects $-18.0 \cdot v_{\text{steer}}$. Item-level paired McNemar counts vs Steering ($+v_{\text{steer}}$): vs BM25 RAG ($b=16, c=60, p=3.85\times 10^{-7}$); vs Hybrid RAG ($b=21, c=27, p=0.4709$). $^\dagger$\,Profiled in separate GPU sessions; total latency reflects shorter completions under retrieved context (decoding 6.81~s) vs full explanations under steering (23.34~s). Retrieval overhead (142.5~ms) measures query encoding, FAISS search, and RRF fusion.
\end{table}

In an operational pipeline comparison, Linear Decay steering (+8.80~pp over BM25 RAG 68.60\%, item-level exact McNemar $p = 3.85 \times 10^{-7}$, paired counts $b=16, c=60$) substantially outperforms sparse lexical retrieval with PyVi tokenization~\cite{ref_pyvi}, which suffers from vocabulary mismatch in Vietnamese clinical terminology (Recall@1 = 19.20\%, MRR@10 = 0.2433). However, against Hybrid RRF RAG (76.20\%, Recall@1 = 48.20\%), steering (77.40\%) yields an observed difference of only +1.20~pp (387 vs 381 wins), which is not statistically significant (paired counts $b=21, c=27$, item-level exact McNemar $p = 0.4709$, paired 95\% bootstrap CI $[-1.60\text{ pp}, +4.00\text{ pp}]$). The operational benefit of early-stopping steering is being entirely retrieval-free: eliminating external passage retrieval indexing, query encoding, and reranking pipelines (142.5~ms retrieval overhead on Kaggle T4) and 215-token context prompt expansion, while operating with an observed in-session runtime difference of +0.01~s (23.34~s vs 23.33~s baseline), unresolved within run-to-run measurement variability ($\pm 0.30$~s).

\subsection{Controlled Steering Strength Ablation}
Initial steering strength ablation (Table~\ref{tab:equal_alpha_ablation}) evaluating finite-window schedules ($K=16$) across initial strengths $\alpha_0 \in \{15.0, 18.0, 20.0\}$ alongside Nominal Matched Dose ($\alpha_{\text{match}} = 0.0425 \alpha_0$) against the continuous reference baseline ($\alpha_0=18.0$) indicates that Linear Decay ($K=16$) achieves peak accuracy at $\alpha_0 = 18.0$ (77.40\% RefPref, 3.88\% Rep-4), serving as an operational evaluation under matched initial scale while acknowledging cumulative dose disparity ($16|\alpha_0|$ for Cutoff vs $8.5|\alpha_0|$ for Decay).

\begin{table}[!t]
\caption{Initial steering strength ablation across intervention schedules ($N_{\text{test}}=500$, $T=200$).}
\label{tab:equal_alpha_ablation}
\centering
\small
\resizebox{\textwidth}{!}{%
\begin{tabular}{llcccc}
\toprule
\textbf{$\alpha_0$} & \textbf{Schedule} & \textbf{RefPref (\%)} & \textbf{BERTScore F1} & \textbf{ROUGE-L (\%)} & \textbf{Rep-4 (\%)} \\
\midrule
-- & Unsteered Baseline & 73.20 & 0.7134 & 23.12 & 3.67 \\
-- & Continuous Reference ($\alpha_0{=}18.0$) & 75.60 & 0.7133 & 23.36 & 3.90 \\
\midrule
\multirow{3}{*}{15.0}
 & Hard Cutoff ($K{=}16$) & \textbf{75.80} & \textbf{0.7146} & 23.11 & 3.73 \\
 & Linear Decay ($K{=}16$) & 75.60 & 0.7142 & 23.10 & 4.03 \\
 & Nominal Matched Dose & 75.40 & 0.7144 & \textbf{23.40} & 3.78 \\
\midrule
\multirow{3}{*}{18.0}
 & Hard Cutoff ($K{=}16$) & 77.00 & 0.7151 & 23.16 & 3.89 \\
 & \textbf{Linear Decay ($K{=}16$)} & \textbf{77.40} & \textbf{0.7157} & 23.33 & \textbf{3.88} \\
 & Nominal Matched Dose & 75.40 & 0.7140 & 23.21 & 4.38 \\
\midrule
\multirow{3}{*}{20.0}
 & Hard Cutoff ($K{=}16$) & 76.40 & 0.7139 & 23.17 & 3.87 \\
 & Linear Decay ($K{=}16$) & \textbf{76.60} & \textbf{0.7147} & 23.15 & \textbf{3.74} \\
 & Nominal Matched Dose & 75.20 & 0.7136 & 23.33 & 4.52 \\
\bottomrule
\end{tabular}%
}
\smallskip\noindent\scriptsize Note: Evaluated on disjoint test split $N_{\text{test}}=500$ ($T=200$). Continuous reference ($K{=}\infty$) is reported once as a fixed global reference baseline evaluated at its optimal scale $\alpha_0=18.0$. Nominal Matched Dose sets per-token coefficient $\alpha_{\text{match}} = 0.0425 \alpha_0$.
\end{table}

\subsection{Temporal Shift, Prefix Disentanglement \& Combined RAG Evaluation}\label{sec:advanced_analysis}
\textbf{Temporal Window Injection Shift Control.} To evaluate whether intervention timing influences factual alignment independently of duration, we compare equal 16-step injection windows under Hard Cutoff ($\alpha_0=18.0$, Layer~8) across matched horizons. Under the controlled 200-token horizon ($T=200$), all 500 completions execute all 199 post-prefill decoding steps without natural EOS emission (EOS hit rate 0.00\%, min 200, max 200 tokens), representing a forced-horizon stress test where every sequence receives the full nominal 16-step dose ($D_i = 16 \times 18.0 = 288$). Factual alignment exhibits monotonic attenuation as injection is deferred: Early intervention ($t \in [1, 16]$) achieves 77.00\% RefPref (+3.80~pp over baseline 73.20\%, item-level exact McNemar $p=0.0110$), Delayed injection ($t \in [17, 32]$) yields 73.40\% (+0.20~pp, 367/500), and Later injection ($t \in [33, 48]$) drops to 72.80\% (364/500; net $-2$ items vs baseline 73.20\%, paired counts $b=16, c=14$, item-level exact McNemar $p=0.8555$). Extended completion ($T=800$) evaluated deferred windows for repetition tracking (5.25\%--5.56\% Rep-4 vs 4.10\% baseline). This monotonic decline (77.00\% $\rightarrow$ 73.40\% $\rightarrow$ 72.80\%) serves strictly as an observational ablation showing attenuated sensitivity when steering is deferred.

\begin{table}[!t]
\centering
\caption{Truncated prefix RefPref accuracy (\%) across sequence length boundaries ($N_{\text{test}}=500$).}
\label{tab:prefix_disentanglement}
\small
\resizebox{\textwidth}{!}{%
\begin{tabular}{lccccc}
\toprule
\textbf{Schedule / Condition} & \textbf{$T{=}50$} & \textbf{$T{=}100$} & \textbf{$T{=}200$} & \textbf{$T{=}400$} & \textbf{$T{=}800$} \\
\midrule
Unsteered Baseline & 71.20 & 72.00 & 73.20 & 67.80 & 65.40 \\
Continuous ($K{=}\infty$) & 74.80 & 75.40 & 75.60 & 72.40 & 68.60 \\
Hard Cutoff ($K{=}16$) & 75.80 & 76.40 & 77.00 & 73.20 & 70.60 \\
\midrule
\textbf{Hard Cutoff Net Gain over Base (pp)} & \textbf{+4.60} & \textbf{+4.40} & \textbf{+3.80} & \textbf{+5.40} & \textbf{+5.20} \\
\bottomrule
\end{tabular}%
}
\smallskip\noindent\scriptsize Note: Evaluated on disjoint test split $N=500$ across sequence truncation boundaries $T \in \{50, 100, 200, 400, 800\}$.
\end{table}

\textbf{Combined Hybrid RAG + Steering Evaluation.} Under natural completion ($T=800$), evaluating combined Hybrid RAG + Steering (Hard Cutoff $K=16, \alpha_0=18.0$) across full $N_{\text{test}}=500$ samples achieves \textbf{80.40\% RefPref accuracy} (+1.00~pp observed difference over standalone Hybrid RAG 79.40\%, paired counts $a=392, b=5, c=10, d=93$, item-level exact binomial McNemar $p=0.3018$, 95\% bootstrap CI $[-0.60\text{ pp}, +2.60\text{ pp}]$). This non-significant increment tests for marginal improvement under combined intervention rather than an interaction effect, which may reflect ceiling effects at ~80\% accuracy, limited statistical power ($N_{\text{discordant}}=15$), or overlapping signals. Combined RAG + Steering evaluates the optimal standalone schedule under extended generation (Hard Cutoff $K=16, \alpha_0=18.0$, 70.60\% vs 67.80\% for Linear Decay) while maintaining low 4-gram repetition (\textbf{4.46\%} vs 4.63\% RAG alone) and concise clinical completions ($167.86$ vs $171.82$ tokens). Selecting Hard Cutoff for the combined RAG experiment was informed post-hoc by its standalone test performance under $T=800$, representing an exploratory joint evaluation on the test split.

\section{Discussion \& Limitations}\label{sec:discussion}
\textbf{Automated Reference Preference vs Clinical Correctness.} Automated BERTScore RefPref evaluates semantic similarity against reference answers, functioning as a preference proxy rather than a clinician-certified metric of medical correctness. Embedding metrics do not perform rule-based symbolic verification of numerical drug dosages, unit accuracy (e.g., mg vs $\mu\text{g}$), clinical contraindications, or drug-drug interactions. Generated model outputs were not subjected to blinded adjudication by practicing physicians or clinical pharmacists; human auditing was restricted to an $N=50$ formulary-grounding spot check on benchmark source pairs (94.0\% alignment, 47/50, 95\% Clopper-Pearson CI $[83.5\%, 98.7\%]$), confirming that ViHaluEval-Medical is a rule-curated benchmark rather than an exhaustively certified clinical ground truth. Additionally, candidate pairs were synthesized using Qwen2.5-7B-Instruct-AWQ, while the evaluated model belongs to the same family (Qwen2.5-7B-Instruct). Although reference preference uses an external encoder (mBERT), this shared architecture introduces potential same-family inductive bias; cross-model generalization (e.g., Llama-3, Gemma-2) remains an essential future test.

\textbf{BERTScore F1 Divergence \& Truncation Sensitivity Analysis.} Under natural completion ($T=800$), pairwise RefPref increases by +5.20~pp while absolute BERTScore F1 against reference answers slightly decreases (0.6695 vs 0.6779) alongside a modest rise in Rep-4 (5.35\% vs 4.10\%). We hypothesize that this divergence relates to output length dynamics (mean 325.6 vs 248.7 Qwen subwords), lowering token-level precision against concise monograph snippets. Regarding evaluator token limits, completions average 276.1--299.2 mBERT subwords; tail completions exceeding mBERT's 512-subword window occur in 7/500 (1.40\%) for Baseline, 6/500 (1.20\%) for Hard Cutoff, 7/500 (1.40\%) for Continuous, and 11/500 (2.20\%) for Linear Decay, undergoing tail truncation in mBERT. In a sensitivity analysis across the 488 mutually non-truncated pairs between Baseline and Hard Cutoff, obtained by componentwise subtraction of the 12 truncated cases ($a_e=7, b_e=0, c_e=1, d_e=4$) from the full table ($311, 16, 42, 131$), relative preference gains remain robust and statistically consistent: Baseline achieves 65.57\% (320/488) while Hard Cutoff achieves 70.70\% (345/488), yielding a +5.12~pp gain ($a=304, b=16, c=41, d=127$, item-level exact McNemar $p=0.0013$, 95\% bootstrap CI $[+2.46, +7.99]$~pp). This sensitivity analysis confirms that evaluator tail truncation beyond 512 subwords (accounting for special tokens) does not drive the observed preference advantage.

\textbf{Intervention Timing vs Cumulative Dose Confounding.} Comparing schedules under matched initial scale $\alpha_0=18.0$ is intrinsically confounded by cumulative intervention dose ($16|\alpha_0|=288$ for Cutoff vs $8.5|\alpha_0|=153$ for Decay vs $L_i|\alpha_0|$, reaching up to $199 \times 18 = 3,582$ for a completion that executes all 199 post-prefill decoding steps). Because intervention duration and cumulative dose co-vary across schedules, current evaluations reflect their joint operational effect; isolating the independent causal role of early timing from cumulative dose requires fine-grained dose-matched ablations across varying horizons. While Nominal Matched Dose ($0.0425\alpha_0$, executing $8.4575\alpha_0$ over $L_{\text{max}}=199$ post-prefill steps) calibrates dose-rate to the nominal 200-token horizon, early EOS emission makes executed dose variable across instances; hence, degradation under Continuous steering cannot be causally isolated from cumulative overdose.

\textbf{Activation Norm Mechanics \& Cluster-Adjusted Uncertainty.} Post-cutoff ($t > K$), within-run offset $\Delta_{\text{norm}} = 0.00$ occurs by design ($\alpha(t)=0$). The cross-condition comparison against unsteered inference ($51.84$ vs $52.20$, paired $t$-test $p=0.1050 > 0.05$, $95\%$ bootstrap CI $[-0.80, +0.08]$) fails to reject equal mean scalar norms, indicating that mean scalar magnitudes do not significantly diverge from baseline without establishing dynamic trajectory equivalence. Test questions cluster by active pharmaceutical ingredient (API) or drug monograph (373 distinct clusters across $N_{\text{test}}=500$, mean size 1.34). While primary tests treat items as item-level observations, intra-cluster correlation across questions sharing a monograph means effective sample size may be smaller than nominal $N=500$, potentially rendering standard item-level McNemar and bootstrap confidence intervals optimistic; cluster-robust variance estimation remains an important direction. Additionally, Holm step-down adjustments were applied locally within individual evaluation panels.

\textbf{Directional Controls \& Deployment Boundaries.} Random label permutation ($N_{\text{flipped}}=5,198/10,272$, 68.60\%) supports the importance of contrastive semantics. Covariance-matched directions ($N=20$, $74.81\% \pm 0.95\%$ sample SD) achieve +1.61~pp over baseline despite near-zero mean cosine similarity with $v_{\text{steer}}$ ($-0.06$). Because zero-mean Gaussian covariance sampling is centrally symmetric ($v \leftrightarrow -v$), sampling from $\Sigma_{\text{data}}$ restricts perturbations to empirical contrastive variance (avoiding isotropic collapse, 55.82\%), while $+v_{\text{steer}}$ supplies directed centroid alignment (+2.59~pp / $2.73\sigma$). We acknowledge that $N=20$ samples and a single permutation split represent preliminary control benchmarks, insufficient to conclusively prove that $v_{\text{steer}}$ isolates an exclusive factual semantic direction. Because vector extraction occurs at sequence-terminal tokens (\texttt{<|im\_end|>}), $v_{\text{steer}}$ may encode stylistic, length, or termination patterns alongside factuality, motivating future token-level pooling. Deployment claims should strictly reflect measured quantities: early-stopping steering provides a lightweight, retrieval-free preference alignment mechanism eliminating external retrieval indexing (142.5~ms) and context prompt expansion with an observed runtime difference of +0.01~s vs baseline (unresolved within run-to-run measurement variability $\pm 0.30$~s), without asserting unmeasured end-to-end clinical deployment guarantees.

\section{Conclusion}\label{sec:conclusion}
We introduced Early-Stopping Activation Steering, a temporally bounded inference-time intervention for Vietnamese clinical Small Language Models. Evaluated on ViHaluEval-Medical (14,674 rule-filtered pairs from 14,700 candidates derived from the Vietnamese National Drug Formulary, 2nd Edition, 2018), restricting steering to initial $K=16$ decoding steps yields reference-preference alignment with zero model updates. Under natural completion ($T=800$), Hard Cutoff ($K=16$) achieves 70.60\% RefPref (+5.20~pp, item-level exact McNemar $p=0.00086$, Holm $p_{\text{adj}}=0.0026$). Under bounded stress testing ($T=200$), Linear Decay ($K=16$) achieves 77.40\% (+4.20~pp, item-level exact McNemar $p=0.0055$, Holm $p_{\text{adj}}=0.0165$), yielding no statistically significant difference from Hybrid RRF RAG (76.20\%, $p=0.4709$) and significantly outperforming lexical BM25 (+8.80~pp, $p=3.85 \times 10^{-7}$) without external retrieval overhead and with minimal runtime difference (+0.01~s vs baseline, within $\pm 0.30$~s measurement variability). While exploratory combination of Hybrid RAG with Hard Cutoff yields 80.40\% (+1.00~pp, $p=0.3018$, non-significant marginal gain), early-stopping steering provides a lightweight, retrieval-free preference alignment mechanism warranting clinician-adjudicated validation.

\begin{credits}
\subsubsection{\ackname} The author acknowledges FPT University Quy Nhon for facilitating computational resources and academic support for this research project.

\subsubsection{\discintname} The author declares no competing financial or non-financial interests directly related to the work described in this paper.
\end{credits}

\begin{thebibliography}{21}
\bibitem{ref1} Zou, A. et al.: Representation engineering: A top-down approach to AI transparency. arXiv:2310.01405 (2023)
\bibitem{ref2} Li, K. et al.: Inference-time intervention: Eliciting truthful answers from a language model. In: Proc. NeurIPS, vol. 36, pp. 41451--41530 (2023)
\bibitem{ref3} Singhal, K. et al.: Large language models encode clinical knowledge. Nature \textbf{620}, 172--180 (2023)
\bibitem{ref4} Qwen Team: Qwen2.5 Technical Report. arXiv:2412.15115 (2024)
\bibitem{ref5} Lewis, P. et al.: Retrieval-augmented generation for knowledge-intensive NLP tasks. In: Proc. NeurIPS, vol. 33, pp. 9459--9474 (2020)
\bibitem{ref6} \vn{Dược thư Quốc gia Việt Nam}: \vn{Ban Dược thư Quốc gia, Nhà xuất bản Y học, Hà Nội} (2018)
\bibitem{ref7} Hu, E.J. et al.: LoRA: Low-rank adaptation of large language models. In: Proc. ICLR (2022)
\bibitem{ref8} Turner, A. et al.: Activation addition: Steering language models without optimization. arXiv:2308.10248 (2023)
\bibitem{ref9} Rimsky, N. et al.: Steering llama 2 via contrastive activation addition. arXiv:2312.06681 (2023)
\bibitem{ref10} Ouyang, L. et al.: Training language models to follow instructions with human feedback. In: Proc. NeurIPS, vol. 35, pp. 27730--27744 (2022)
\bibitem{ref11} Azaria, A., Mitchell, T.: The internal state of an LLM knows when it's lying. In: Proc. EMNLP, pp. 967--976 (2023)
\bibitem{ref12} Kirkpatrick, J. et al.: Overcoming catastrophic forgetting in neural networks. Proc. PNAS \textbf{114}(13), 3521--3526 (2017)
\bibitem{ref13} Li, J. et al.: HaluEval: A large-scale hallucination evaluation benchmark for large language models. In: Proc. EMNLP, pp. 6449--6464 (2023)
\bibitem{ref14} Lin, C.-Y.: ROUGE: A package for automatic evaluation of summaries. In: Text Summarization Branches Out, pp. 74--81 (2004)
\bibitem{ref15} Zhang, T. et al.: BERTScore: Evaluating text generation with BERT. In: Proc. ICLR (2020)
\bibitem{ref16} Lin, S. et al.: TruthfulQA: Measuring how models mimic human falsehoods. In: Proc. ACL, pp. 3214--3252 (2022)
\bibitem{ref_vllm} Kwon, W. et al.: Efficient memory management for large language model serving with PagedAttention. In: Proc. SOSP, pp. 611--626 (2023)
\bibitem{ref_bnb} Dettmers, T. et al.: QLoRA: Efficient finetuning of quantized LLMs. In: Proc. NeurIPS, vol. 36, pp. 10088--10115 (2023)
\bibitem{ref_bgem3} Chen, J. et al.: BGE M3-Embedding: Multi-lingual, multi-functionality, multi-granularity text embeddings through self-knowledge distillation. arXiv:2402.03216 (2024)
\bibitem{ref_faiss} Johnson, J. et al.: Billion-scale similarity search with GPUs. IEEE Trans. Big Data \textbf{7}(3), 535--547 (2019)
\bibitem{ref_pyvi} Nguyen, V.-H.: PyVi: Python Vietnamese core NLP toolkit. Python package (2016)
\end{thebibliography}

\end{document}
